\documentclass[11pt]{article}
\usepackage{fontspec}
\setmainfont{Times New Roman}
\setsansfont{Arial}
\setmonofont{Consolas}
\usepackage{amsmath,amssymb,mathtools}
\usepackage{geometry}
\usepackage{microtype}
\newcommand{\mo}{\mathfrak{o}}
\newcommand{\mT}{\mathfrak{T}}
\newcommand{\mP}{\mathfrak{P}}
\newcommand{\mpideal}{\mathfrak{p}}
\geometry{a4paper,margin=2.35cm}
\setlength{\parindent}{1.25em}
\setlength{\parskip}{0pt}
\makeatletter
\renewcommand\footnotesize{\@setfontsize\footnotesize{7.7}{8.7}\selectfont}
\makeatother
\emergencystretch=100em
\allowdisplaybreaks
\sloppy
\hbadness=10000
\vbadness=10000
\hfuzz=2pt
\vfuzz=10pt
\raggedbottom

\begin{document}

% Унит: Папер31 Сецтион08 Ентры07 дисцриминант идеал ундер пассаге то qуотиент ринг в001. Дирецт Интерславиц Латин транслатион фром цлеан РА34 лине 441 / сегмент P31-С0100, вит сцан-висибле Х. Грелл прооф фоотноте ресторед фром принтед паге 103 / ПДФ паге 22.
\setcounter{footnote}{37}

\noindent 4. \emph{Дискриминантны идеал при преходу к колцу квоциентов.} Нехај $\mpideal$ буде просты идеал из $\mo$, и означмо чрез $\mP=\mT\mpideal$ разширјены идеал в $\mT$. Колцо квоциентов $\mT_\mP$ имаје модулну базу составјену из линеарно независных елементов взходно к колцу главных идеалов $\mo_\mpideal$; Једночасно $\mT_\mP$ јест конечны $\mpideal$-порјадок в разширјеном польу $K$ польа квоциентов колца $\mo_\mpideal$.\footnote{За доказ сравн. Х. Грелл.} Тым самым в $\mT_\mP$ взходно к $\mo_\mpideal$ јест исполньен дискриминантны теорем (§7, 3).

\end{document}

